\section{Finite Calibration and Protocol Comparison}
\label{sec:estimation}

Calibration measurements constrain covariance directions that are invisible
under the realised protocol. We consider $m$ independent units observed densely
according to
\begin{equation}
W^{(i)}=Z^{(i)}+\eta^{(i)},\qquad
Z^{(i)}\sim\cN(0,K),\qquad \eta^{(i)}\sim\cN(0,R_0),
\qquad i=1,\dots,m.
\label{eq:calibration}
\end{equation}
The pairs $(Z^{(i)},\eta^{(i)})$ are independent across $i$, and
$\eta^{(i)}$ is independent of $Z^{(i)}$ for every $i$. Consequently,
$\Cov(W^{(i)})=K+R_0$.
The finite-sample question is how error in the covariance recovered from these
data propagates to comparisons among protocols. Appendix~\ref{sec:app-calibration}
gives the proofs, including covariance-repair and resolvent constants.

\subsection{Covariance Calibration}

Let $\widehat\Sigma=m^{-1}\sum_iW^{(i)}W^{(i)\top}$, and set
$\widehat R_0=R_0$ when the calibration-noise covariance is known;
otherwise estimate $R_0$ from an independent or replicated calibration
component. For a symmetric matrix $M$, define
\begin{equation}
\begin{aligned}
\widetilde M_\tau
&=\proj_{\{H=H^\top:\,H\succeq\tau I_p\}}(M),\\
D_\tau(M)&=\Diag\{\diag(\widetilde M_\tau)\},&
\cR_\tau(M)&=D_\tau(M)^{-1/2}\widetilde M_\tau D_\tau(M)^{-1/2}.
\end{aligned}
\label{eq:repair-operator}
\end{equation}
Thus $\cR_\tau$ floors the eigenvalues and then rescales the result to a
correlation matrix. The plug-in estimator is
\begin{equation}
\widehat K=\cR_{\tau_m}(\widehat\Sigma-\widehat R_0),
\qquad
\widehat\cI_g(S)=\frac{F_g(S;\widehat K)}{V_g(\widehat K)},
\label{eq:plugin}
\end{equation}
where $\tau_m\downarrow0$. The projection in
\eqref{eq:repair-operator} is the Frobenius-norm eigenvalue-clipping map
\citep{higham1988computing}. The asymptotic results require only
$\tau_m\downarrow0$.

The target enters the error rate through the modulus of continuity of its
Gaussian covariance transform $C_g$.

\begin{proposition}[Regularity of the target covariance transform]
\label{prop:regularity}
For each of the following cases there is a finite constant $L$ such that, on
the stated range,
\begin{equation}
|C_g(r_2)-C_g(r_1)|\le L|r_2-r_1|^\beta.
\label{eq:target-modulus}
\end{equation}
\begin{enumerate}[label=(\roman*),leftmargin=2.2em,itemsep=2pt]
\item\label{lem:lipschitz} If $g$ belongs to the first-order Gaussian Sobolev
space $W^{1,2}(\phi)$, then
\eqref{eq:target-modulus} holds on $[-1,1]$ with $\beta=1$.
\item\label{lem:holder} If $g_c(z)=\Ind\{z>c\}$, then it holds on $[-1,1]$
with $\beta=1/2$. This exponent cannot be improved uniformly, because
\begin{equation}
C_{g_c}(1)-C_{g_c}(1-\delta)
\sim \frac{e^{-c^2/2}}{2\pi}\sqrt{2\delta},
\qquad \delta\downarrow0.
\label{eq:boundary-asymptotic}
\end{equation}
\item\label{lem:local-lipschitz} For the same threshold target and every
$0\le r_{\max}<1$, \eqref{eq:target-modulus} holds on
$[-r_{\max},r_{\max}]$ with $\beta=1$.
\end{enumerate}
\end{proposition}

Smooth targets therefore transmit covariance error linearly. Threshold targets
do so linearly at a fixed model whose relevant correlations are separated from
$\pm1$, but only through a square-root envelope when boundary correlations are
allowed.

\subsection{Uniform Value Error}

We now propagate covariance-calibration error through the protocol-value
calculation,
$\widehat K-K\mapsto Q_S(\widehat K)-Q_S(K)\mapsto
\widehat\cI_g(S)-\cI_g(S)$. Two quantities can destabilise this chain and
therefore appear explicitly in the assumptions below: a vanishing target
variance and nearly singular protocol covariances.

\begin{theorem}[Uniform stability of protocol values]
\label{thm:uniform-error}
Let $V_g(K)\ge v_0>0$, and suppose that for a feasible family $\Pi_{\mathcal B}$,
\begin{equation}
\sup_{\substack{S\in\Pi_{\mathcal B}\\S\ne\varnothing}}
\opnorm{A_S}^2\le a<\infty,
\qquad
\inf_{\substack{S\in\Pi_{\mathcal B}\\S\ne\varnothing}}
\lambda_{\min}(A_SKA_S^\top+R_S)\ge\lambda>0.
\label{eq:family-conditioning}
\end{equation}
If the modulus \eqref{eq:target-modulus} with exponent $\beta$ is valid for all
arguments of $C_g$ arising from entries of $K$, $\widehat K$, $Q_S(K)$ and
$Q_S(\widehat K)$, then, for all sufficiently small
$e=\opnorm{\widehat K-K}$,
\begin{equation}
\sup_{S\in\Pi_{\mathcal B}}
\big|\widehat\cI_g(S)-\cI_g(S)\big|\le C e^\beta.
\label{eq:uniform-error}
\end{equation}
Here $C$ depends only on the target modulus, $v_0$, bounds on $K$ and
$\widehat K$, and the family bounds $(a,\lambda)$, not on $S$.
Consequently $\beta=1$ for smooth targets, $\beta=1/2$ globally for threshold
targets, and $\beta=1$ for threshold targets at an interior model.
\end{theorem}

For a threshold target, the interior condition invoked in the theorem can be
checked at the fixed model through
\begin{equation}
r_0=\max\left\{\max_{j\ne k}|K_{jk}|,\ 
\sup_{S\in\Pi_{\mathcal B}}\max_{j,k}|Q_S(K)_{jk}|\right\}<1.
\label{eq:r0}
\end{equation}
Uniform continuity of the resolvent then places the corresponding arguments
for every sufficiently close $\widehat K$ in a common compact subset of
$(-1,1)$.

\subsection{Calibration Rates}

\begin{theorem}[Fixed-model calibration rates]
\label{thm:fixed-model}
Fix $p$ and a finite feasible family $\Pi_{\mathcal B}$, and suppose that
$K\succ0$, $V_g(K)>0$ and the conditioning bound in
\eqref{eq:family-conditioning} holds. If
$\opnorm{\widehat R_0-R_0}=O_p(m^{-1/2})$ and $\tau_m\downarrow0$, then the
estimator in \eqref{eq:plugin} satisfies
$\opnorm{\widehat K-K}=O_p(m^{-1/2})$. Hence
\begin{equation}
\sup_{S\in\Pi_{\mathcal B}}|\widehat\cI_g(S)-\cI_g(S)|=
\begin{cases}
O_p(m^{-1/2}), & g\in W^{1,2}(\phi),\\
O_p(m^{-1/2}), & g=g_c\text{ and the interior condition \eqref{eq:r0} holds},\\
O_p(m^{-1/4}), & g=g_c\text{ under the global threshold envelope}.
\end{cases}
\label{eq:root-m}
\end{equation}
\end{theorem}

The global $m^{-1/4}$ rate is a worst-case guarantee over correlation matrices;
the ordinary root-$m$ rate applies to every fixed threshold model satisfying
\eqref{eq:r0}.

\subsection{Selection Regret and Distinguishable Value Gaps}

Uniform value error translates directly into a decision guarantee.

\begin{corollary}[Selection regret and distinguishable value gaps]
\label{cor:regret}
Let
$\varepsilon_m=\sup_{S\in\Pi_{\mathcal B}}|\widehat\cI_g(S)-\cI_g(S)|$, let $S^*$
maximise $\cI_g$ over $\Pi_{\mathcal B}$, and let $\widetilde S$ be a feasible
protocol with empirical optimisation gap
\begin{equation}
\eta_m=\max_{S\in\Pi_{\mathcal B}}\widehat\cI_g(S)
-\widehat\cI_g(\widetilde S).
\label{eq:opt-gap}
\end{equation}
Then
\begin{equation}
\cI_g(S^*)-\cI_g(\widetilde S)\le2\varepsilon_m+\eta_m.
\label{eq:regret-opt}
\end{equation}
Exact empirical maximisation is the special case $\eta_m=0$.

For nested classes
$\Pi^{(1)}\subseteq\cdots\subseteq\Pi^{(L)}=\Pi_{\mathcal B}$, let
$\widehat S_\ell$ maximise $\widehat\cI_g$ over $\Pi^{(\ell)}$. On the event
$\sup_{S\in\Pi^{(\ell)}}|\widehat\cI_g(S)-\cI_g(S)|\le\varepsilon_\ell$,
\begin{equation}
\cI_g(S^*)-\cI_g(\widehat S_\ell)\le
\underbrace{\cI_g(S^*)-
\max_{S\in\Pi^{(\ell)}}\cI_g(S)}_{\text{class restriction}}
+\underbrace{2\varepsilon_\ell}_{\text{calibration error}}.
\label{eq:resolution-bound}
\end{equation}
\end{corollary}

An estimated gap larger than $2\varepsilon_m$ certifies the same ordering at
the population level, whereas a population gap larger than $2\varepsilon_m$
cannot be reversed by the plug-in estimates. Hence protocol-value differences
below this scale cannot be distinguished uniformly from calibration error.
Equation~\eqref{eq:resolution-bound} separates this uncertainty from the
approximation loss incurred by restricting the search to $\Pi^{(\ell)}$.

\theoremstyle{remark}
\newtheorem*{bestlinearremark}{Remark}
\begin{bestlinearremark}[Best-linear calibration]
The same perturbation and regret arguments apply to $\cI_{\mathrm L}$. If the
protocol covariances are uniformly nonsingular and their covariance and
target-cross-covariance blocks, together with the target variance, are
estimated uniformly at the root-$m$ rate, then
$\sup_{S\in\Pi_{\mathcal B}}|\widehat\cI_{\mathrm L}(S)-\cI_{\mathrm L}(S)|
=O_p(m^{-1/2})$. Appendix~\ref{sec:app-best-linear} states sufficient
conditions and proves this claim; \cref{cor:regret} then applies with
$\cI_g$ replaced by $\cI_{\mathrm L}$.
\end{bestlinearremark}
